% Legendas LaTeX das tres figuras de AIME (Latour 2013), Etapa 3.
% Arquivos PNG em outputs/latour_2013_aime_en/figuras/. Apos copiar os
% PNG para o diretorio de figuras da tese, ajuste o path em
% \includegraphics. O \autoref{sec:aime-tres-camadas} na figura 2 remete
% a uma secao do capitulo 2 que pode ser rotulada conforme a estrutura
% que adotar.
% Estilo das legendas: primeira pessoa do singular, sem economia de
% ressalvas, sem desenvolver interpretacao que exceda o que a contagem
% mostra. As pistas que aparecem nas legendas (registro deslocado do
% militar, articulacao topologia+trajetoria_passe, distribuicao
% homogenea de \texttt{Moderns} ao longo do livro) estao todas
% documentadas no \texttt{outputs/etapa3_aime/relatorio_etapa3.md}.

% =====================================================================
% Figura 1: frequencia dos 26 campos figurativos com ocorrencias.
% =====================================================================

\begin{figure}[htbp]
\centering
\includegraphics[width=\linewidth]{figuras/aime_frequencia_grupos.png}
\caption[Frequencia dos campos figurativos em \emph{AIME}]{
Frequencia absoluta dos vinte e seis campos figurativos com
ocorrencias em \emph{An Inquiry into Modes of Existence} (Latour,
2013), o livro com que Latour declara estar completando a teoria
ator-rede (capitulo 12: \enquote*{the actor-network theory, which the
present inquiry seeks to complete}). A coloracao separa os campos por
proveniencia do catalogo lexical do meu projeto. Em azul, os campos do
catalogo estabelecido nas Etapas 1 e 2 da analise computacional, que
rastreiam o vocabulario teorico de Latour entre 1986 e 1999: dos
dezenove campos do catalogo, quatorze aparecem em \emph{AIME}, com
destaque para \texttt{topologia} (514 ocorrencias) e \texttt{network}
(255), em densidades comparaveis as dos livros monograficos anteriores.
Em laranja, os doze campos novos do catalogo especifico que construi
para esta Etapa 3, identificados a partir da leitura qualitativa do
livro e da literatura sobre a guinada latouriana posterior a 2007:
\texttt{trajetoria\_passe} (433 ocorrencias) e \texttt{modernos} (397)
lideram, sinalizando que o vocabulario figurativo novo que \emph{AIME}
introduz e tao denso quanto o vocabulario antigo persistente. A
operacao teorica do livro nao e deslocamento total nem continuidade
simples: e justaposicao. Soma de 3.557 ocorrencias em 194.454 palavras
de corpo. Sem validacao amostral semantica nesta etapa; os campos de
alta polissemia (\texttt{experiencia}, \texttt{categorias\_dominio})
carregam ruido nao estimado, conforme registro em
\texttt{outputs/etapa3\_aime/relatorio\_etapa3.md} secao 8. Figura
gerada por \texttt{scripts/23\_etapa3\_aime\_visualizacoes.py} a partir
de \texttt{outputs/latour\_2013\_aime\_en/csv/kwic\_catalogo\_antigo.csv}
e \texttt{kwic\_catalogo\_aime.csv}.}
\label{fig:aime-frequencia-grupos}
\end{figure}


% =====================================================================
% Figura 2: densidade dos 12 campos mais frequentes ao longo do texto.
% =====================================================================

\begin{figure}[htbp]
\centering
\includegraphics[width=\linewidth]{figuras/aime_densidade_top12.png}
\caption[Densidade dos doze campos mais frequentes ao longo de
\emph{AIME}]{Distribuicao das ocorrencias dos doze campos figurativos
mais frequentes ao longo de \emph{An Inquiry into Modes of Existence}
(Latour, 2013), em histograma empilhado de trinta bins. O eixo
horizontal representa a posicao relativa de cada ocorrencia no texto
(zero no inicio, um no fim). A leitura visual mostra tres coisas que
importam ao meu argumento sobre a trajetoria latouriana 1986--2013. A
primeira: \texttt{modernos} e \texttt{Moderns} (em laranja-claro,
laranja-medio na paleta) aparecem em densidade alta e de modo
relativamente homogeneo ao longo de todo o livro, sustentando a
designacao central da \enquote{antropologia dos modernos} que da
subtitulo a obra. A segunda: \texttt{topologia} e
\texttt{trajetoria\_passe} concentram densidade muito alta na primeira
metade do livro, onde Latour expoe o aparato conceitual dos modos de
existencia. A terceira: \texttt{militar} (decimo na densidade
absoluta, com 129 ocorrencias) aparece distribuido ao longo do texto
sem concentracao em capitulo especifico, em contraste com o pico
inicial que esse mesmo campo apresentou em \emph{Science in Action}
1987 e em \emph{Pandora's Hope} 1999, e em registro deslocado para o
horizonte cosmopolitico-diplomatico, conforme detalhamento da secao 5
do relatorio da Etapa 3. A alternancia entre campos do catalogo antigo
(\texttt{topologia}, \texttt{network}, \texttt{textil},
\texttt{militar}) e campos do catalogo novo da Etapa 3
(\texttt{trajetoria\_passe}, \texttt{modernos},
\texttt{categorias\_dominio}, \texttt{experiencia},
\texttt{instituicao}, \texttt{modos\_existencia}, \texttt{alteracao},
\texttt{felicidade}) reforca a leitura de \emph{AIME} como livro de
tres camadas figurativas: vocabulario antigo persistente, vocabulario
novo especifico, e articulacao entre os dois evidenciada pela
cocorrencia (vide \autoref{sec:aime-tres-camadas}). Paleta tab20 do
matplotlib. Fonte: \texttt{kwic\_catalogo\_antigo.csv} e
\texttt{kwic\_catalogo\_aime.csv} (coluna \texttt{posicao\_no\_texto}
em offsets de caractere). Figura gerada por
\texttt{scripts/23\_etapa3\_aime\_visualizacoes.py}.}
\label{fig:aime-densidade-top12}
\end{figure}


% =====================================================================
% Figura 3: densidade de todos os 26 campos ao longo do texto.
% =====================================================================

\begin{figure}[htbp]
\centering
\includegraphics[width=\linewidth]{figuras/aime_densidade_todos.png}
\caption[Densidade de todos os campos figurativos ao longo de
\emph{AIME}]{Distribuicao das ocorrencias dos vinte e seis campos
figurativos com presenca em \emph{An Inquiry into Modes of Existence}
(Latour, 2013), em histograma empilhado de trinta bins, com posicao
relativa no eixo horizontal (zero no inicio, um no fim). Versao
auditavel da \autoref{fig:aime-densidade-top12}, em que mantenho todos
os campos com ocorrencia maior que zero (incluindo os menos densos:
\texttt{diplomacia} com 84 ocorrencias, \texttt{preposicao} com 71,
\texttt{gaia} com 52, \texttt{construction} com 74, e demais).
Importante para a leitura: cinco campos do catalogo antigo das Etapas
1 e 2 zeram em \emph{AIME} (\texttt{centre\_of\_calculation},
\texttt{trial\_of\_strength}, \texttt{factish},
\texttt{circulating\_reference}, \texttt{spokesperson}) e por isso nao
aparecem aqui. O \texttt{factish} e o \texttt{circulating\_reference},
em particular, sao vocabulario caracteristico de \emph{Pandora's Hope}
1999 (4,14 e 1,41 por dez mil palavras) que \emph{AIME} abandona; o
\texttt{trial\_of\_strength}, vocabulario agonistico de
\emph{Science in Action} 1987, tambem desaparece. Esse abandono
silencioso de vocabulario anterior, somado a reorganizacao das
variantes top do \texttt{topologia} (que em \emph{AIME} sao
\texttt{path} e \texttt{trajectory}, articuladas com
\texttt{trajetoria\_passe}), e parte do que a Etapa 3 documenta sobre
o trabalho que \emph{AIME} faz com seu proprio passado. Fonte:
\texttt{kwic\_catalogo\_antigo.csv} e \texttt{kwic\_catalogo\_aime.csv}.
Figura gerada por
\texttt{scripts/23\_etapa3\_aime\_visualizacoes.py}.}
\label{fig:aime-densidade-todos}
\end{figure}
